\documentclass{beamer}

% for adding conditionals
\usepackage{etoolbox}

% for adding note support
\usepackage{pgfpages}
% only compile notes if `\def\notes{}` in preamble
\ifdef{\notes}{\setbeameroption{show notes on second screen}}{}

\usetheme{CambridgeUS}
\usecolortheme{dolphin}
\usefonttheme{serif}

\definecolor{rit@orange}{HTML}{F36E21}
\definecolor{rit@brown}{HTML}{513127}
\setbeamercolor{structure}{fg=rit@orange}
\setbeamercolor{palette primary}{fg=white, bg=rit@brown}
\setbeamercolor{palette secondary}{fg=rit@brown}
\setbeamercolor{palette tertiary}{fg=white, bg=rit@orange}
\setbeamercolor{titlelike}{fg=rit@brown}
\setbeamercolor{normal text}{fg=rit@brown}

\setbeamertemplate{enumerate}[circle]
\setbeamertemplate{items}[circle]

\AtBeginSection{\frame{\sectionpage}}

\setbeamertemplate{section page}
{
    \begin{centering}
    \begin{beamercolorbox}[sep=12pt,center]{part title}
    \usebeamerfont{section title}\insertsection\par
    \end{beamercolorbox}
    \end{centering}
}

% remove shadow behind blocks
\setbeamertemplate{blocks}[rounded][shadow=false]
% set blocks to use custom color scheme
\setbeamercolor{block title}{bg=rit@orange,fg=white}
\setbeamercolor{block body}{bg=black!5,fg=rit@brown}
% remove shading between block title and body
\makeatletter
\pgfdeclareverticalshading[lower.bg,upper.bg]{bmb@transition}{200cm}{%
  color(0pt)=(lower.bg); color(4pt)=(lower.bg); color(4pt)=(upper.bg)}
\makeatother

\setbeamercovered{dynamic}
\beamertemplatetransparentcoveredmedium

\beamertemplatenavigationsymbolsempty

%\logo{\vspace{-0.25cm}\includegraphics[width=1in]{img/tiger_walking_rit_color}}
\logo{%
  \makebox[0.95\paperwidth]{%
    \includegraphics[width=.75in,keepaspectratio]{img/tiger_walking_rit_color}%
    \hfill%
    \includegraphics[width=.75in,keepaspectratio]{img/lucifiers-screwdriver-logo}%
  }%
}

\pdfpageattr {/Group << /S /Transparency /I true /CS /DeviceRGB>>}

\usepackage{amsmath,amssymb,commath,mathabx,physics,siunitx}
\usepackage{cclicenses,copyrightbox}
\usepackage{float,subcaption}





\let\svthefootnote\thefootnote
\textheight 1in
\newcommand\blankfootnote[1]{%
  \let\thefootnote\relax\footnotetext{#1}%
  \let\thefootnote\svthefootnote%
}


\title{Lucifer's Screwdriver}

\author{Daniel Wysocki}

\institute[RIT]{
  Rochester Institute of Technology
}


\date[04/24/2016]{
  NASA Space Apps Challenge
  \\
  April 23-24, 2016
}

\begin{document}

\maketitle

\begin{frame}{Challenge}

  \begin{itemize}
  \item Develop a:
    \begin{itemize}
    \item hypothetical method
    \item concept note
    \item simple prototype
    \end{itemize}
  \item demonstrating how machine learning can help us avoid the same
    fate as the dinosaurs.
  \end{itemize}

\end{frame}


\begin{frame}{Solution}

  \begin{center}
    $k$-Nearest Neighbor classifier to identify

    \vspace{0.25in}

    \textbf{Potentially Hazardous Objects (PHA)}
  \end{center}

\end{frame}


\begin{frame}{Potentially Hazardous Objects (PHA)}

  \begin{itemize}
  \item Two important criteria
    \begin{enumerate}
    \item Will it hit us?
    \item Is it big?
    \end{enumerate}
  \item We quantify this with two things
    \begin{enumerate}
    \item Minimum Orbit Intersection Distance (MOID) less than
      $\SI{0.05}{AU}$.
    \item Infrared brightness exceeds a minimum threshold ($H < 22$).
    \end{enumerate}
  \end{itemize}

\end{frame}


\begin{frame}{MOID Calculation}

  \begin{itemize}
  \item Take orbital parameters measured by the Minor Planet Center
    \begin{itemize}
    \item $a$, $e$, $i$, $\omega$, $\Omega$
    \end{itemize}
  \item Use existing simulation software to calculate the MOID between
    the Earth and each object.
  \end{itemize}

\end{frame}


\begin{frame}{$k$-NN Classification}

{
  \Large
  \begin{displaymath}
    (a, e, i, \omega, \Omega, H) \overset{k\mathrm{NN}}{\to} \mathrm{PHA?}
  \end{displaymath}
}

\end{frame}


\begin{frame}{Lucifer's Screwdriver}

  \begin{figure}[ht]
    \centering
    \includegraphics[width=0.4\textwidth]{img/lucifiers-screwdriver-logo}

    \url{https://github.com/dwysocki/lucifers-screwdriver/}
  \end{figure}

  \vfill

\end{frame}



\end{document}
